
\newpage
\setcounter{table}{0}
\setcounter{figure}{0}
\section{绪论}

\subsection{研究背景与意义}

专利文献及其结构化元数据是技术创新活动的重要外显载体。与普通科技文本相比，专利数据同时具有"技术属性、法律属性、时间属性、组织属性"四重特征，这使其天然适合用于技术演化分析、竞争态势分析和创新能力评估\cite{basberg1987patents,trajtenberg1990penny,daim2006forecasting}。长期以来，专利计量分析在科技政策制定、企业技术战略规划和创新能力评估中发挥着不可替代的作用\cite{griliches1990patent,narin1997increasing,zhu2014patent}。专利引用数据作为知识流动的重要指标，已被广泛用于测度技术影响力和知识溢出效应\cite{hall2001nber,jaffe2000knowledge}。

然而，在实际研究场景中，专利计量分析仍高度依赖人工流程：研究者通常需要手工定义变量、挑选统计方法、设置控制变量并反复调整分析路径。该过程的专业门槛高、耗时长、复现实验成本大，且在面对新问题时迁移效率较低\cite{vanzeebroeck2011filing,mejia2024patent}。

近年来，大语言模型（Large Language Model, LLM）推动了自动化数据分析系统的快速发展\cite{zhao2024llm,brown2020language}。已有工作在任务分解、代码生成、执行纠错等方面取得显著进展，例如 DS-Agent\cite{guo2024dsagent}、Data Interpreter\cite{hong2025data} 等系统展示了多智能体协作在复杂分析任务中的可行性。多智能体框架如 AutoGen\cite{wu2023autogen}、MetaGPT\cite{hong2024metagpt} 等进一步提供了可编程的协作机制，使得构建复杂的端到端分析系统成为可能\cite{wang2024survey}。

但现有系统的优化重点主要集中在"执行层"，即如何更稳定地完成既定分析任务；对于"规划层"，即如何在给定研究问题与具体数据集时自动形成高质量分析方案，现有方法仍不充分\cite{sun2025survey}。该问题在专利计量场景中尤为突出，因为专利变量关系复杂、字段语义多样、跨列关联隐含，若缺乏前置数据感知，系统易产生模板化、泛化但不贴合数据事实的方案。

另一方面，知识图谱增强 LLM 已成为热点方向。GraphRAG\cite{edge2024graphrag} 等研究表明，图结构能够提升复杂问题的全局信息组织能力。在知识图谱与 LLM 的融合路线上，已有系统性的路线图综述\cite{pan2024unifying}。在专利领域，EvoPat\cite{wang2024evopat}、KLIPA\cite{deng2025klipa} 等工作开始尝试将 LLM 与知识组织或检索能力结合。

但总体上，已有研究更多服务于问答、摘要、检索等文本任务；针对"专利结构化数据上的分析方案自动设计"仍缺少系统性方法\cite{jiang2025natural}。由此，本文聚焦"数据感知驱动的多智能体专利分析方案生成"，试图补足从执行自动化到规划自动化的关键环节。

从实践价值看，若能在规划前自动提取数据分布、图结构、样本语义等证据信息，并将其与领域知识约束协同使用，可显著降低人工试错成本，提高研究可复现性与结论可解释性。对专利计量研究者、技术情报分析师和企业知识产权部门而言，这一能力均具有直接应用价值。

\subsection{研究问题与研究内容}

围绕上述背景，本文的核心研究问题定义为：如何构建一个面向专利计量分析场景的多智能体系统，使其能够在给定具体专利数据集与研究问题的前提下，自动生成具有数据感知能力的分析方案，并在执行反馈基础上持续迭代优化，从而稳定提升分析结果的针对性、可解释性与整体质量。与仅关注代码生成或单步任务执行的自动化方法不同，本文更关注“分析方案如何被规划出来”这一前置问题，即如何让系统在正式执行之前，先形成与当前数据特征、研究目标和领域知识相匹配的分析思路。

为回答该问题，本文围绕方案生成过程中的“数据依据从何而来”“理论与方法如何约束”“系统如何持续改进”三个关键环节，开展以下三项研究内容：

第一，提出数据感知方案生成机制。系统在正式规划前，并不直接根据自然语言问题生成分析步骤，而是首先通过 DataPreview 提取列级统计、字段类型、缺失情况与跨列关系等信息，通过 GraphPreview 提取专利网络中的结构特征，并结合分层抽样摘要触发 LLM 生成“可检验洞察”。这些洞察并非泛化表述，而是能够被后续统计分析、建模过程或代码执行直接验证的候选判断。该方法的目标，是将“数据事实”显式注入规划上下文，使后续方案生成不再只依赖语言模型对任务表述的语义理解，而是建立在对当前数据集真实结构的感知基础之上，从而减少模板化、空泛化和脱离样本实际的问题。

第二，设计双知识图谱协同机制。本文将因果图谱用于约束“研究假设空间”，将方法图谱用于约束“分析执行空间”，并在规划阶段实现二者联动\cite{pan2024unifying,liu2016knowledge}。其中，因果图谱主要回答“哪些变量关系值得被优先考察、哪些路径具备理论上的可解释性”，方法图谱则进一步回答“这些变量可以如何被测度、这些关系适合采用什么分析方法进行验证”。通过将两类图谱分别作用于假设形成与方法选择两个层面，系统能够在保持一定生成灵活性的同时，避免分析路径完全依赖大模型自由发挥。该机制的核心价值，在于把专利分析中的领域知识、变量结构与方法经验沉淀为可调用、可组合、可约束的图谱资源，从而提升方案的理论一致性与执行可行性。

第三，构建两轮迭代优化机制。第一轮聚焦探索性分析，目标是在数据感知结果与图谱约束的基础上快速形成较为完整的初始分析蓝图，并识别出具有潜在价值的变量关系、比较维度和方法路径；第二轮则基于第一轮执行结果进行差距识别、方案修正和验证性深化，形成“探索—验证”闭环。与一次性输出方案的方式相比，两轮迭代设计能够让系统显式吸收执行阶段暴露出的不足，例如变量选择不充分、方法与数据不匹配、结论支撑证据不足等问题，并据此对后续规划进行补强。该方法有助于提升分析深度、机制解释能力和结论稳健性，使系统输出更接近真实研究过程中的渐进式推理与修正逻辑。

\subsection{研究方法与技术路线}

本文采用“系统构建 + 渐进叠加 + 组件消融 + 重复实验 + 双轨评估”的研究方法。整体上，本文并非仅从算法性能提升的角度展开研究，而是遵循“先构建可运行系统原型，再通过渐进实验和消融实验识别关键模块贡献，最后以多维指标和重复实验验证方法有效性”的技术路线推进。该路线兼顾了系统研究的工程实现要求与学术研究的可解释性要求，能够较为完整地回答“系统是否可实现”“改进是否真实有效”“提升究竟来自何种机制”以及“这些提升是否在重复运行下仍然成立”等核心问题。

在系统构建层面，本文基于工作流编排框架\cite{langgraph2024}设计四智能体协作流程：规划智能体负责分析蓝图生成，方法规格化智能体负责任务技术规格化，执行智能体负责代码生成与纠错执行，评审智能体负责结果校验与报告生成。系统以状态机方式组织任务流，支持结构化输入、过程追踪与结果落盘。在这一基础上，本文进一步将数据预览、图结构预览、双知识图谱调用、任务规格化与结果反馈等能力嵌入统一工作流中，使系统不只是“多代理串联”，而是形成从问题理解、方案规划、代码执行到结果回评的闭环式分析流程。这样的设计有助于提升系统在复杂专利计量任务中的模块协同性、过程透明度与故障可定位性。

在实验层面，本文采用四方案渐进叠加设计：D\_baseline0（纯 LLM）、A\_template（仅知识图谱模板）、B\_data\_aware（知识图谱+数据洞察）、C\_iterative（知识图谱+数据洞察+两轮迭代，完整模式）。该设计遵循“由简到繁、逐层加能”的思想，通过控制变量的方式依次引入知识约束、数据感知和迭代优化机制，使每一个模块的边际贡献都能够被较清晰地观察和解释。与直接比较“完整系统”和“单一基线”的做法相比，这种渐进叠加实验更有利于揭示系统性能变化背后的原因，避免出现“整体提升但原因不明”的评估陷阱，同时也为后续分析知识图谱、数据洞察和两轮迭代之间的协同关系提供了实验基础。

在评估层面，本文联合两类指标：其一为自动化结构指标（列覆盖率、方法多样性、参数丰富度、任务数等）；其二为 LLM-as-Judge 语义指标（相关性、针对性、方法合理性、创新性、深度），该方法已在 LLM 评估领域得到广泛验证\cite{zheng2023judging}。前者强调系统输出在结构层面的完整性、丰富性与可复核性，适合衡量方案是否真正覆盖了数据分析任务所需的关键组成部分；后者则更关注最终回答研究问题时的语义质量，能够从更接近人类评审的角度评估方案是否“有用、合理、深入”。在此基础上，本文进一步对 6 种模式在 3 个问题上各独立重复运行 5 次，共形成 90 组实验结果，以降低单次观察对总体结论的影响。

\subsection{研究难点与创新点}

本文研究存在两个核心难点。与一般的自动化文本处理任务不同，本文所面向的是“面向具体专利数据集的分析方案自动生成”问题，其难点不仅体现在语言模型是否能够理解问题语义，更体现在系统能否将数据事实、领域知识与分析执行过程有机整合。因此，难点并非单纯来自模型能力不足，而是来自规划约束、方法匹配与迭代反馈之间的系统性耦合。

第一，如何让分析方案与数据事实保持高一致性。现实中 LLM 容易输出“表面合理”的泛化路径，但未必对应当前数据集的真实分布与结构特征。尤其在专利计量场景中，不同数据集在字段构成、时间跨度、样本规模、网络结构和语义密度上往往存在显著差异，若系统不能在规划前充分识别这些差异，就容易生成看似完整、实则脱离数据基础的分析蓝图。若无法将数据证据转化为规划约束，自动化系统将难以超越提示词工程式模板，也难以在复杂研究问题上形成真正具有针对性的方案。

第二，如何平衡“知识约束”与“生成自由度”。知识图谱可以提升结构一致性，帮助系统在变量关系、方法路径和研究逻辑上避免明显失真，但过强约束也可能压缩模型的探索空间，使系统倾向于重复已有模板，而难以生成更具启发性的分析组合。特别是在专利分析这类交叉性较强的问题上，研究者往往既希望方案“有依据”，又希望方案“有新意”。因此，如何找到约束与开放之间的可操作平衡，使知识图谱既成为支撑规划质量的“护栏”，又不沦为限制创新能力的“天花板”，是系统设计中的关键问题。

针对上述难点，本文形成三点创新。这些创新并非彼此孤立，而是分别对应“数据从何而来”“知识如何作用”“系统为何能够持续变好”三个核心问题，共同支撑本文提出的多智能体专利分析方案生成框架。

\begin{enumerate}[(1)]
\item 提出“数据感知前置”的方案生成流程。通过 DataPreview、GraphPreview 与洞察生成模块，将数据证据前移至规划阶段，形成“数据先证据化、方案后生成”的机制。该创新的关键不在于单纯增加几个统计特征，而在于改变了分析方案的生成顺序，使系统先理解“当前数据能支持什么”，再决定“应该做什么分析”，从而显著增强方案与数据基础之间的贴合性。
\item 提出双知识图谱协同架构。将“因果解释”与“方法执行”拆分到两个图谱并在规划阶段联动，提升方案的理论一致性与执行可行性。相较于将所有知识混合到同一模板中的做法，这种分层式图谱设计更有利于分别约束研究假设空间与方法选择空间，使系统既能够在理论上保持结构清晰，又能够在执行上形成可落地的方法路径。
\item 通过重复实验识别知识约束的稳定收益边界。90 组重复实验表明，知识图谱并非在所有问题上都以同样强度发挥作用：在结构锚点清晰、分析链条较长的问题上，图谱能够稳定提升方案组织与方法选择；在开放度更高的问题上，静态路径约束则需要与数据证据更紧密协同，才能避免压缩模型的探索空间。该发现将知识图谱的作用从“是否有效”进一步细化为“在什么条件下更有效”，为后续构建自适应约束强度机制提供了实验依据。
\end{enumerate}

\subsection{论文结构安排}

全文共七章，各章之间按照“问题提出与文献梳理一系统设计与实现一实验验证与应用展示一总结与展望”的逻辑展开，既体现从理论分析到系统落地的递进关系，也体现从方法提出到效果验证的完整研究闭环。各章内容安排如下。

第一章为绪论，主要说明本文的研究背景、问题定义、研究方法、研究难点与创新点，并对全文结构进行总体说明。通过本章的论述，明确本文的研究对象、研究目标以及所要解决的核心问题，为后续章节展开奠定总体框架。

第二章为研究综述，系统梳理 LLM 自动化分析、多智能体协作框架、LLM 在专利分析中的应用，以及知识图谱增强 LLM 的研究进展，并在此基础上归纳现有工作的主要不足。通过文献综述，本文进一步明确自身的研究切入点，即面向专利结构化数据场景，补足自动化分析系统在“规划层”能力上的缺口。

第三章为系统设计方案，围绕本文提出的多智能体专利分析系统，给出设计目标、总体架构、核心能力层以及典型应用场景映射。本章重点回答“系统为什么这样设计”“各模块之间如何协同”以及“数据感知、双知识图谱与两轮迭代机制如何共同嵌入整体架构”等问题，是全文方法设计层面的核心章节。

第四章为系统原型实现，进一步从工程实现角度描述系统关键模块的落地方式，包括数据预览、图结构预览、数据洞察生成、双知识图谱调用、分析蓝图生成与执行闭环，以及两轮迭代机制的具体实现过程。本章重点体现系统从设计思想走向可运行原型的实现路径，展示各模块的输入输出关系与技术实现细节。

第五章为实验与评估，介绍实验数据集、实验方案设计、评估指标体系以及不同模式下的结果对比与机制分析。本章通过渐进叠加实验、组件消融实验和多维评价指标，对系统的有效性、稳定性与改进来源进行系统验证，从而回答“系统是否真正优于基线方案”以及“提升主要来自哪些关键机制”等问题。

第六章为系统应用示例，通过多个完整案例展示系统从输入研究问题、数据感知、方案规划、代码执行到结果生成的端到端流程。本章重点体现系统在真实专利分析任务中的应用方式与输出效果，使前文提出的设计与实验结论在具体场景中得到进一步印证，也增强全文的应用导向与实践说服力。

第七章为总结与展望，对全文研究工作进行归纳总结，提炼本文在方法、系统与实验层面的主要结论，并结合当前研究仍存在的局限讨论未来可能的拓展方向。通过本章，进一步说明本文工作的研究价值及其在后续自动化专利分析与知识增强智能系统中的延伸潜力。
